\documentclass[10pt,a4paper]{amsart}
\usepackage[margin=19mm]{geometry}
\usepackage{amsmath,amssymb,mathtools}
\usepackage[scr=rsfs,cal=euler]{mathalfa}
\usepackage{xfrac}
\usepackage[unicode,hidelinks,bookmarks=false]{hyperref}
\newcommand{\ie}{i.e.\ }
\newcommand{\cf}{cf.\ }
\newcommand{\resp}{respectively\ }
\newcommand{\Subsection}{\subsection}
\providecommand{\nobreakdash}{\nobreak\hbox{-}}
\setlength{\emergencystretch}{2em}
\multlinegap0pt
\usepackage{amssymb}
\usepackage{enumerate}
\usepackage{stackengine}
\usepackage{mathtools}
\usepackage{tikz}
\newtheorem{theorem}{Theorem}
\title{On the approximation of Weierstrass function via superoscillations}
\author{Fabrizio Colombo, Irene Sabadini, Daniele C. Struppa}
\date{Source: arXiv:2603.05580v1 (CC BY 4.0)}
\begin{document}
\maketitle

\noindent\emph{Source and licence.} On the approximation of Weierstrass function via superoscillations. Version arXiv:2603.05580v1 (CC BY 4.0), distributed by arXiv under the Creative Commons Attribution 4.0 International licence, http://creativecommons.org/licenses/by/4.0/, as recorded in the arXiv licence metadata for this exact version and shown on the abstract page https://arxiv.org/abs/2603.05580v1. Full licence notice: \texttt{licenses/arxiv-2603.05580-CC-BY-4.0.txt}.\par\smallskip

\noindent\textit{Editorial scope.} Counted unit: the article's theorem thm:carnot (source0.tex lines 232--241). The article's complete original proof of it is delivered with it (source0.tex lines 243--302, one \texttt{proof} environment of 60 lines). No mathematics is written, repaired or re-derived: the statement and the proof are the article's own lines, in the article's own notation, and no retained line is substituted or re-flowed.  No further source-local context is needed: every cross-reference of every retained line resolves inside the two delivered spans.  Nothing was omitted from any retained span, so the package carries no omission record.  No retained line contains a citation, so the wrapper carries no bibliography and no bibliography entry.  No retained line contains a figure environment, an image-inclusion command or drawing code, so no image is delivered and the package declares no asset dependency.  Editorial formatting only, in addition: an AMS \texttt{amsart} preamble that restates the article's own declarations and macros used by the retained lines, namely the environments \texttt{theorem} on separate counters, together with the standard packages those lines need.  The retained environments print in this document as Theorem 1 in order of appearance, whereas the article prints the same items with the numbering of its own full document.  The complete native source of the article as distributed in the arXiv e-print for this exact version is bundled unchanged as \texttt{sources/195812/source0.tex}.
\begin{theorem}[Explicit error estimate] \label{thm:carnot}
Let $M > 0$ be a fixed real number and $\alpha \in \mathbb{R}$. Define the constants:
\[
    K(\alpha) = \frac{|\alpha^2-1|M^2}{2} \quad \text{and} \quad J(\alpha) = 2|\alpha(1-\alpha^2)|M^3.
\]
Then, for any integer $n \ge \frac{4M}{\pi}$, the following explicit estimate holds for all $|x| \le M$:
\begin{equation}
    \left| \left( \cos\frac{x}{n} + i \alpha \sin\frac{x}{n} \right)^n - e^{i \alpha x} \right| \leq \frac{1}{n} \sqrt{ K(\alpha)^2 e^{2K(\alpha)/n} + \frac{J(\alpha)^2}{n^2} e^{K(\alpha)/n} }.
\end{equation}
\end{theorem}

\begin{proof}
Let $w := F_n(x, \alpha) = \left(\cos\frac{x}{n} + i\alpha \sin\frac{x}{n}\right)^n$ and $z := e^{i\alpha x}$. We aim to bound the Euclidean distance $|w-z|$ in the complex plane. Representing both numbers in polar form, $w = \rho_w e^{i\theta_w}$ and $z = \rho_z e^{i\theta_z}$, we have:
\begin{align*}
    \rho_w &= \left(\cos^2 \frac{x}{n} + \alpha^2 \sin^2 \frac{x}{n}\right)^{n/2} = \left( 1 + (\alpha^2-1)\sin^2 \frac{x}{n} \right)^{n/2}, \\
    \theta_w &= n \arctan \left(\alpha \tan \frac{x}{n}\right), \\
    \rho_z &= 1, \quad \theta_z = \alpha x.
\end{align*}
By the Carnot's Theorem in the complex plane:
\begin{equation} \label{eq:carnot_id}
    |w-z|^2 = \rho_w^2 + 1 - 2\rho_w \cos(\theta_w - \theta_z) = (\rho_w - 1)^2 + 2\rho_w(1 - \cos(\theta_w - \theta_z)).
\end{equation}
Using the standard Taylor inequality $1 - \cos y \leq \frac{y^2}{2}$, valid for all $y \in \mathbb{R}$, we obtain:
\begin{equation} \label{eq:main_ineq_proof}
    |w-z|^2 \leq (\rho_w - 1)^2 + \rho_w (\theta_w - \theta_z)^2.
\end{equation}

\medskip
We estimate of the modulus $\rho_w$
using the inequality $|(1+y)^\gamma - 1| \leq \gamma |y| e^{\gamma |y|}$ for $\gamma > 0$ and $y > -1$. Setting $\gamma = n/2$ and $y = (\alpha^2-1)\sin^2(x/n)$, and noting that $|\sin u| \le |u|$, we find:
\[
    |\rho_w - 1| \leq \frac{n}{2} |\alpha^2-1| \frac{x^2}{n^2} \exp\left( \frac{n}{2} |\alpha^2-1| \frac{x^2}{n^2} \right) \leq \frac{|\alpha^2-1|M^2}{2n} \exp\left( \frac{|\alpha^2-1|M^2}{2n} \right).
\]
By substituting $K(\alpha) = \frac{|\alpha^2-1|M^2}{2}$, we obtain:
\begin{equation} \label{eq:rho_est}
    |\rho_w - 1| \leq \frac{K(\alpha)}{n} e^{K(\alpha)/n}.
\end{equation}
Furthermore, since $\rho_w = (\rho_w - 1) + 1 \le |\rho_w - 1| + 1$, it follows from the growth of the exponential that $\rho_w \leq e^{K(\alpha)/n}$.

\medskip
Now we estimate of the phase difference $\Delta \theta = \theta_w - \theta_z$
define
$$h(u) = n \arctan(\alpha \tan u) - n \alpha u,
$$
 such that $\Delta \theta = h(x/n)$. Let $g(u) = \arctan(\alpha \tan u) - \alpha u$. Its derivative is:
\[
    g'(u) = \frac{1}{1 + (\alpha \tan u)^2} \cdot \alpha \sec^2 u - \alpha = \alpha \left( \frac{\sec^2 u - (1 + \alpha^2 \tan^2 u)}{1 + \alpha^2 \tan^2 u} \right) = \frac{\alpha(1-\alpha^2)\tan^2 u}{1+\alpha^2 \tan^2 u}.
\]
By the Mean Value Theorem, there exists $\xi$ between $0$ and $x/n$ such that
$$|g(x/n) - g(0)| = |x/n| |g'(\xi)|.
$$
 Since $n \ge 4M/\pi$, we have $|\xi| < |x/n| \le \pi/4$. In this range, $\tan^2 \xi \le 2\xi^2$ (since $\tan \xi \approx \xi$ and $1+\alpha^2 \tan^2 \xi \ge 1$). Thus:
\[
    |g'(\xi)| \leq |\alpha(1-\alpha^2)| \tan^2 \xi \leq 2 |\alpha(1-\alpha^2)| \xi^2 \leq 2 |\alpha(1-\alpha^2)| \frac{x^2}{n^2}.
\]
Then
$$|\Delta \theta| = n |g(x/n)| \leq n \cdot \frac{|x|}{n} \cdot 2 |\alpha(1-\alpha^2)| \frac{x^2}{n^2} \leq \frac{2 |\alpha(1-\alpha^2)| M^3}{n^2}.
$$
Setting $J(\alpha) = 2|\alpha(1-\alpha^2)|M^3$, we get $|\Delta \theta| \leq \frac{J(\alpha)}{n^2}$.

\medskip
Substituting the estimates into \eqref{eq:main_ineq_proof}:
\[
    |w-z|^2 \leq \left( \frac{K(\alpha)}{n} e^{K(\alpha)/n} \right)^2 + e^{K(\alpha)/n} \left( \frac{J(\alpha)}{n^2} \right)^2.
\]
Factoring out $1/n^2$ from the entire expression:
\[
    |w-z|^2 \leq \frac{1}{n^2} \left[ K(\alpha)^2 e^{2K(\alpha)/n} + \frac{J(\alpha)^2}{n^2} e^{K(\alpha)/n} \right].
\]
Taking the square root of both sides yields the desired inequality.
\end{proof}

\end{document}
